\chapter{Discovery, Investigation, and Master Chronology}
\label{chap:discovery-investigation-master-chronology}

\classification{Evidence status: Chronology reconstruction. Confidence baseline: High for source-mapped verified events; Moderate where events are partially verified; Unknown for later unverified modern-status entries retained only as chronology leads. This chapter does not evaluate hypotheses or infer motive, cause, or identity beyond the recorded chronology.}

\section*{Chronology Control}

This chapter reconstructs the case chronology from the master chronology database. It admits only events recorded in the authorized master chronology, with uncertainty preserved from the chronology uncertainty register and source links preserved from the event-source map. The purpose is chronological orientation, not explanation.

The chronology contains 24 admitted events. Eleven are verified, 11 are partially verified, and 2 later-period modern-status entries are unable to verify from the acquired primary-source archive. Eight unsupported or unable-to-verify timeline claims were removed or logged as uncertainties rather than smoothed into the sequence.

Dates and times are presented only at the precision supported by the source. Unknown times remain unknown. Witness-reported times remain witness-reported. Approximate medical estimates remain approximate. A chronological gap is not filled by inference.

\begin{table}[htbp]
\centering
\caption{Chronology verification summary}
\label{tab:chronology-verification-summary}
\begin{tabularx}{\textwidth}{>{\RaggedRight\arraybackslash}p{0.34\textwidth}r>{\RaggedRight\arraybackslash}X}
\toprule
Chronology category & Count & Use in this chapter \\
\midrule
Verified events & 11 & Used as the primary factual backbone. \\
Partially verified events & 11 & Used with explicit limits or uncertainty language. \\
Unable-to-verify modern-status entries & 2 & Included only as modern-status markers requiring official documentation. \\
Removed or unsupported timeline claims & 8 & Excluded from narrative sequence and retained only as uncertainty controls. \\
\bottomrule
\end{tabularx}
\end{table}

\section{Pre-Discovery Context}

No verified or partially verified pre-1948 event is admitted into the master chronology. The current evidence base does not support a pre-discovery biography, travel history, intelligence context, or identity chronology for the unknown man at the level required for this chapter.

The first admitted chronology entries begin on 30 November 1948 with transport-ticket evidence. The Henley Beach railway ticket and the Adelaide-Somerton bus ticket are important because they are associated with property found on or with the body. They do not, by themselves, identify the passenger or reconstruct the full route taken.

On 30 November 1948, ticket clerk Douglas George Townsend was on duty at the Adelaide Railway Station ticket box from 6:15 a.m. to about 2:00 p.m. and later stated that the Henley Beach ticket was issued during his shift. The associated second-class Adelaide-to-Henley Beach ticket was issued that day, but the exact issue time and purchaser identity remain unknown.

The bus-ticket evidence is more specific about route timing but not passenger identity. The chronology records that a bus ticket associated with the body was issued on the Adelaide-Somerton route departing the Railway Station at about 11:15 a.m. The holder, boarding point, and alighting point remain uncertain.

\section{Discovery at Somerton Beach}

The evening witness sequence begins at 7:00 p.m. on 30 November 1948. A man was reported lying where the body was later found, in a similar position, and was seen to raise his right arm before it fell limply. The event is verified in the chronology, but the identity of the man seen was not proven by facial identification.

At about 7:30 p.m., a man and a girl saw a man in the same place. They did not see movement, and one witness formed an impression that he was lying unnaturally. This event is also verified as a witness observation. It does not establish with certainty that the person seen was the deceased.

The master chronology records a negative-evidence interval between the bus arrival and the 7:00 p.m. sighting: the coroner recorded that no one came forward to say the man was seen at Somerton during that interval. This is a statement about the received witness record, not proof that no one saw him.

On 1 December 1948, at about 7:00 a.m., the body of an unknown man was found on the shore at Somerton Beach. This is the central chronology anchor and is classified as verified with very high confidence. Exact first-discovery witness details still require page-level witness-statement confirmation.

Dr Bennett examined the body at about 9:40 a.m. and formed the opinion that death had occurred about eight hours earlier. The chronology also records a medical estimate placing death around about 2:00 a.m. Both entries are partially verified because they are medical timing estimates, not exact observed times of death.

\section{Initial Police Investigation}

The early police sequence is documented but not complete in time detail. The body was taken to the City Mortuary, but the exact transfer time is not established in the extracted chronology.

An initial search of the clothing found a Henley Beach railway ticket, a bus ticket, cigarettes, combs, and chewing gum. The Tamam Shud slip was not found in this initial search. The exact time of the search remains unknown.

The witness index records Constable John Moss as a first responding police witness who described attendance after a telephone message, body position, a cigarette near the collar or cheek, lack of visible violence, transport, and initial property search. The chronology uses these details only where they are mapped to admitted events and source pages.

Identification attempts began in the early investigation but are not reconstructed beyond the admitted chronology. Fingerprints, photographs, public circulation, and missing-person comparisons are relevant investigative tools, but the master chronology does not admit a solved identification event for 1948 or 1949.

\section{Post-Mortem Investigation}

On 2 December 1948, Robert James Cowan received specimens including stomach and contents, liver and muscle, urine, and blood from the unidentified man found at Somerton the previous day. This event is verified with high confidence.

Cowan's analysis is recorded separately. The chronology states that he was unable to find signs of any common poison in the submitted specimens. This is partially verified with high confidence because it supports negative common-poison testing, not the absence of all possible poisons.

The medical and post-mortem chronology should therefore be read narrowly. It establishes specimen receipt and negative common-poison testing in the recorded evidence. It does not establish exact cause of death, a specific poison, or absence of all toxic agents.

\section{Expansion of the Investigation}

At a later date recorded only as December 1948 or later, the Tamam Shud slip was found in a trouser fob pocket that was difficult to locate. The event is verified with high confidence, but the exact discovery date, discovery time, and finder identity remain unresolved in the chronology.

The Tamam Shud paper was associated by the coroner with a copy of the second edition of Fitzgerald's translation of the Rubaiyat of Omar Khayyam. This event is partially verified with high confidence. The book-recovery chronology and provenance remain unresolved.

A little over a month after 1 December 1948, an unclaimed suitcase was found at the Adelaide Railway Station luggage office. This event is partially verified with high confidence. The exact date, property records, and chain of custody remain gaps.

The suitcase was assessed as having internal evidence connecting it with the deceased or his death, including clothing size, removed tags, and thread similarity. This is a partially verified chronology event. It records the coroner's assessment of connection; it does not establish complete ownership, complete inventory, or complete custody.

Public appeals and interstate enquiries appear in the witness and source layers, but the master chronology admits them only where source-mapped event records exist. In this chapter, identification efforts are therefore described as part of the investigative expansion, not as resolved identity findings.

\begin{table}[htbp]
\centering
\caption{Major investigative milestones}
\label{tab:chronology-investigative-milestones}
\begin{tabularx}{\textwidth}{>{\RaggedRight\arraybackslash}p{0.18\textwidth}>{\RaggedRight\arraybackslash}p{0.22\textwidth}>{\RaggedRight\arraybackslash}X >{\RaggedRight\arraybackslash}p{0.18\textwidth}}
\toprule
Event ID & Date or period & Milestone & Status \\
\midrule
CHR-0009 & 1948-12-01 & Body found on the shore at Somerton Beach. & Verified, Very High \\
CHR-0011 & 1948-12-01 & Initial clothing search found transport tickets and small property; Tamam Shud slip not found. & Partially Verified, High \\
CHR-0012 & 1948-12-01 & Body taken to the City Mortuary. & Verified, High \\
CHR-0013 & 1948-12-02 & Cowan received post-mortem specimens. & Verified, High \\
CHR-0014 & 1948-12-02 & Cowan found no signs of common poison in submitted specimens. & Partially Verified, High \\
CHR-0015 & December 1948 or later & Tamam Shud slip found in a difficult-to-locate trouser fob pocket. & Verified, High \\
CHR-0017 & A little over a month after 1948-12-01 & Unclaimed suitcase found at Adelaide Railway Station luggage office. & Partially Verified, High \\
CHR-0019 & 1949-06-17 & 1949 inquest proceedings documented in DOC-0001. & Verified, High \\
\bottomrule
\end{tabularx}
\end{table}

\section{Inquest Timeline}

The 1949 inquest period begins in the admitted chronology on 17 June 1949. The 1949 inquest file includes proceedings and remarks concerning the death of the body located at Somerton. The event is verified with high confidence from the primary-document inventory and \texttt{DOC-0001} page references.

The inquest chronology runs from 17 June to 21 June 1949. The coroner adjourned the inquest sine die after recording that the deceased was unknown, that death was not natural, and that precise circumstances and cause could not be stated with certainty from the reviewed evidence. This records the inquest status; it is not converted here into a modern cause-of-death conclusion.

The witness index identifies medical, police, transport, railway, beach, forensic, and exhibit witnesses in the 1949 inquest reproduction. Their testimony is used in this chronology only where it maps to admitted chronology events. Full witness reliability assessment belongs in the later testimony chapter.

\section{Later Developments}

The later chronology is intentionally limited. A later inquest or coronial file is identified as GRG 1/27 File 53/1958 and dated 14 March 1958. This event is partially verified because the document was acquired as a later reproduction but content requires visual review and original-custodian access remains restricted.

Two archival processing events are admitted after 1958. The preserved 1949 inquest reproduction includes an administrative return or date-stamp sequence indicating archival handling before digitization. A certified copy of the unknown man's death certificate was issued by the Registrar on 14 May 2014. These events concern documentation and preservation history, not the original death event.

The 2021 exhumation is retained only as an unable-to-verify modern-status entry in the Mission 5 chronology. Source-register entries list 2021 reporting and an official-release pathway, but the relevant SAPOL or Forensic Science SA primary release was not acquired and reviewed in the primary-source archive used for Mission 4.

The 2022-2026 Webb identity and genetic-genealogy chronology is also retained only as an unable-to-verify modern-status entry in this chronology layer. The project records the researcher-led Carl Charles Webb attribution elsewhere with official-status caveats. This chapter does not treat it as an officially verified identity event because no acquired official DNA, coroner, SAPOL, or Forensic Science SA report closes that issue in the reviewed corpus.

\section*{Master Chronology Table}

The master timeline is reproduced below from the Mission 5 LaTeX table. It preserves event IDs, dates, times, locations, event summaries, source references, and verification status.

\begingroup
\scriptsize
\setlength{\tabcolsep}{2pt}
\hbadness=10000
\hfuzz=40pt
% Auto-generated Mission 5 timeline table. Do not edit manually without updating research/chronology/chronology_master.csv.
\begin{longtable}{p{0.12\textwidth}p{0.14\textwidth}p{0.13\textwidth}p{0.27\textwidth}p{0.15\textwidth}p{0.11\textwidth}}
\caption{Mission 5 Source-Mapped Master Timeline}\label{tab:mission5-master-timeline}\\
\toprule
Event & Date/Time & Location & Event & Source & Status \\
\midrule
\endfirsthead
\toprule
Event & Date/Time & Location & Event & Source & Status \\
\midrule
\endhead
CHR-0002 & 1948-11-30 / 06:15-14:00 & Adelaide Railway Station ticket box & Ticket clerk Douglas George Townsend was on duty and later stated that the Henley Beach ticket was issued during his shift. & DOC-0001 (p.75) & Verified (High) \\
CHR-0003 & 1948-11-30 / Unknown within 06:15-14:00 & Adelaide Railway Station ticket box & A second-class Adelaide-to-Henley Beach railway ticket associated with the body was issued. & DOC-0001 (p.75) & Verified (High) \\
CHR-0004 & 1948-11-30 / 11:15 a.m. & Adelaide-Somerton bus route from Railway Station toward St Leonards & A bus ticket associated with the body was issued on the Adelaide-Somerton bus route departing the Railway Station at about 11:15 a.m. & DOC-0001 (pp.5,23) & Verified (High) \\
CHR-0005 & 1948-11-30 / Between bus arrival and 7:00 p.m. & Somerton Beach area & The coroner recorded that no one came forward to say the man was seen at Somerton between arrival of the bus and 7:00 p.m. & DOC-0001 (p.7) & Partially Verified (Moderate) \\
CHR-0006 & 1948-11-30 / 7:00 p.m. & Somerton Beach, where body was later found & A man was reported lying where the body was later found, in a similar position, and was seen to raise his right arm before it fell limply. & DOC-0001 (p.5) & Verified (High) \\
CHR-0007 & 1948-11-30 / About 7:30 p.m. & Somerton Beach, where body was later found & A man and a girl saw the man in the same place; they did not see movement and one formed an impression that he was lying unnaturally. & DOC-0001 (p.5) & Verified (High) \\
CHR-0008 & 1948-12-01 / About 2:00 a.m. & Unknown & Dr Bennett estimated death occurred around about 2:00 a.m. & DOC-0001 (p.5) & Partially Verified (Moderate) \\
CHR-0009 & 1948-12-01 / About 7:00 a.m. & Somerton Beach & The body of an unknown man was found on the shore at Somerton. & DOC-0001; DOC-0003 (DOC-0001 pp.5,7) & Verified (Very High) \\
CHR-0010 & 1948-12-01 / 9:40 a.m. & Somerton Beach / body location & Dr Bennett examined the body and formed the opinion that death had occurred about eight hours previously. & DOC-0001 (p.5) & Partially Verified (Moderate) \\
CHR-0011 & 1948-12-01 / Unknown & Body/clothing search; later City Mortuary & A search of the clothing found a Henley Beach railway ticket, bus ticket, cigarettes, combs and chewing gum; the slip bearing Tamam Shud was not found in this initial search. & DOC-0001 (pp.5,21) & Partially Verified (High) \\
CHR-0012 & 1948-12-01 / Unknown & City Mortuary & The body was taken to the City Mortuary. & DOC-0001 (p.21) & Verified (High) \\
CHR-0013 & 1948-12-02 / Unknown & Postmortem / analytical chain; exact location not established & Robert James Cowan received specimens including stomach and contents, liver and muscle, urine and blood from the unidentified man found at Somerton the previous day. & DOC-0001 (p.47) & Verified (High) \\
CHR-0014 & 1948-12-02 / Unknown & Analytical laboratory; exact location not established & Cowan carried out analyses of the specimens and was unable to find signs of any common poison. & DOC-0001 (p.47) & Partially Verified (High) \\
CHR-0015 & December 1948 or later / Unknown & Clothing/body examination; exact location not established & The Tamam Shud slip was found in a trouser fob pocket that was difficult to locate. & DOC-0001 (pp.5,81) & Verified (High) \\
CHR-0016 & December 1948 or later / Unknown & Investigation records; exact location not established & The Tamam Shud paper was associated by the coroner with a copy of the second edition of Fitzgerald's translation of the Rubaiyat of Omar Khayyam. & DOC-0001 (p.5) & Partially Verified (High) \\
CHR-0017 & A little over a month after 1948-12-01 / Unknown & Adelaide Railway Station luggage office & An unclaimed suitcase was found at the Adelaide Railway Station luggage office. & DOC-0001 (pp.5,7) & Partially Verified (High) \\
CHR-0018 & A little over a month after 1948-12-01 / Unknown & Adelaide Railway Station luggage office / property examination & The suitcase was assessed as having internal evidence connecting it with the deceased or his death, including clothing size, removed tags and thread similarity. & DOC-0001 (p.7) & Partially Verified (High) \\
CHR-0019 & 1949-06-17 / Unknown & Coroner proceedings, Adelaide & The 1949 inquest file includes proceedings and remarks concerning the death of the body located at Somerton. & DOC-0001 (DOC-0001 metadata; pp.5-9,37-47,75,81,85) & Verified (High) \\
CHR-0020 & 1949-06-17 to 1949-06-21 / Unknown & Coroner proceedings, Adelaide & The coroner adjourned the inquest sine die after recording that the deceased was unknown, death was not natural, and precise circumstances/cause could not be stated with certainty. & DOC-0001 (pp.7,9) & Verified (High) \\
CHR-0021 & 1958-03-14 / Unknown & Coronial record / State Records SA reference & A later inquest/coronial file is identified as GRG 1/27 File 53/1958. & DOC-0002; SRC-0002 (DOC-0002 metadata; Mission 3 inventory) & Partially Verified (Moderate) \\
CHR-0022 & 2014-05-14 / Unknown & Births, Deaths and Marriages Registration Office, Adelaide & A certified copy of the unknown man death certificate was issued by the Registrar. & DOC-0003 (DOC-0003 p.1) & Partially Verified (High) \\
CHR-0024 & 2009-04-15 / Unknown & State Coroner / archive return process & The preserved 1949 inquest reproduction includes an administrative return/date stamp sequence indicating archival handling before digitization. & DOC-0001 (p.1 OCR/front matter) & Partially Verified (Moderate) \\
CHR-0025 & 2021-05-19 / Morning pending & West Terrace Cemetery & Reported exhumation of remains for modern forensic/DNA work remains outside the acquired primary-source archive reviewed in Mission 4. & SRC-0008; SRC-0009; SRC-0010 listed in source register but not acquired primary archive & Unable to Verify (Unknown) \\
CHR-0026 & 2022-07-26 to 2022-07-27 / Unknown & Public/research announcement and SAPOL response context & Modern claim identifying the man as Carl/Charles Webb remains unsupported by the acquired primary-source archive. & SRC-0011 listed as pending source; no acquired DNA report & Unable to Verify (Unknown) \\
\bottomrule
\end{longtable}

\endgroup

\section{Chronology Assessment}

The strongest chronology anchors are the discovery of the body at Somerton Beach on 1 December 1948, the associated transport-ticket evidence, the 30 November evening witness observations, the initial property search, the City Mortuary transfer, the 2 December specimen receipt, the negative common-poison analysis, the later Tamam Shud slip discovery, the suitcase lead, and the 1949 inquest proceedings.

The partially verified chronology is equally important because it prevents overstatement. The medical estimate of death time remains approximate. The initial search time is unknown. The Tamam Shud discovery date is unresolved. The Rubaiyat book recovery and provenance remain incomplete. The suitcase date, inventory, and custody remain incomplete. The 1958 file is known as a record but not fully reviewed in this chronology.

The largest chronology gaps concern the purchaser and exact issue time of the Henley Beach ticket, the holder and precise use of the bus ticket, whether the man seen on the evening of 30 November was the deceased, the exact time and place of death, the timing of clothing searches, the recovery sequence for the Tamam Shud slip and Rubaiyat copy, the suitcase custody trail, and the official status of modern exhumation and identity developments.

\begin{table}[htbp]
\centering
\caption{Chronology uncertainty summary}
\label{tab:chronology-uncertainty-summary}
\begin{tabularx}{\textwidth}{>{\RaggedRight\arraybackslash}p{0.28\textwidth}>{\RaggedRight\arraybackslash}X >{\RaggedRight\arraybackslash}p{0.22\textwidth}}
\toprule
Uncertainty area & Principal unresolved issue & Related records \\
\midrule
Transport tickets & Purchaser, holder, and exact route use are not established. & CHR-0002 to CHR-0004 \\
Evening sightings & Witness observations do not prove identity continuity. & CHR-0006; CHR-0007 \\
Death timing and place & Medical timing is approximate and place of death remains unresolved. & CHR-0008; CHR-0010 \\
Initial search & Exact search time is unknown and the Tamam Shud slip was found later. & CHR-0011; CHR-0015 \\
Rubaiyat material & Book recovery chronology and provenance remain unresolved. & CHR-0016 \\
Suitcase & Exact discovery date, inventory, ownership, and custody remain incomplete. & CHR-0017; CHR-0018 \\
Later official status & 2021 exhumation and Webb identity developments require official primary documentation. & CHR-0025; CHR-0026 \\
\bottomrule
\end{tabularx}
\end{table}

The chronology supports a disciplined reconstruction of events but not a completed explanation. It establishes what the evidence currently allows the manuscript to place in sequence. It also marks where the sequence stops.
